\chapter{Modules}\label{app-modules}

The Revised Report has no modules: a program is a single text, and the
only library is the standard one. Large programs need a way to divide
the text into separately compiled parts with defined interfaces. The
GNU dialect of Algol 68, which ga68 and a68j implement, adds modules for
this purpose, following a proposal by Lindsey and Boom~\cite{modules};
outside this appendix, the book uses them only for \code|Genie_Transput|. This
appendix describes the essentials. Algol 68 Genie has no modules.

\section*{Defining a module}

A module is a separate source file of the form
\begin{lstlisting}
module Name =
def
   declarations and expressions
fed
\end{lstlisting}
The declarations between \code|def| and \code|fed| are private, unless
they are marked with \code|pub|; the public ones form the interface of
the module. Types, constants, variables, functions, operators and
precedences can be public. The sequence is executed once, before the
program that uses the module starts, and must end with an expression, as
every sequence must (often \code|skip|).

The following module implements a stack of integers. The list behind it
is private; only the three functions are public:

\program{appE/stack}

\section*{Using a module}

A program uses modules by naming them after \code|access|, before its
block:

\example{appE/use-stack}

The public names of all accessed modules are visible in the block, as if
they had been declared in a block around it. Two modules that both
declare the same public name cannot be accessed together, since the name
would be declared twice; a module that is used by two other modules can
be accessed by both, and is executed only once. A module can itself
access other modules, by writing \code|access| before its \code|def|.

\section*{Compiling}

A module is compiled on its own, and the program is compiled with the
compiled modules. With ga68, a module is compiled to an object file,
which also contains its interface, and the program is linked with it:
\begin{lstlisting}[style=shell]
$ ga68 -c stack.a68
$ ga68 -I. use-stack.a68 stack.o genietransput.o -o use-stack
\end{lstlisting}
With a68j, a module is compiled to a library jar, and the directory of
the jar is given to the program:
\begin{lstlisting}[style=shell]
$ java -jar a68j-backend.jar stack.a68 -d lib/libstack.jar
$ java -jar a68j-backend.jar -I lib -run use-stack.a68
\end{lstlisting}
The modules that come with a68j, such as \code|Genie_Transput|, need no
such step.
